\documentclass[12pt,a4paper]{amsart}

\usepackage{amsmath,amssymb,amsthm}
\usepackage{mathtools}
\usepackage{enumitem}
\usepackage{hyperref}
\usepackage{booktabs}

%%─── Theorem environments ────────────────────────────────────────────
\newtheorem{theorem}{Theorem}[section]
\newtheorem{lemma}[theorem]{Lemma}
\newtheorem{proposition}[theorem]{Proposition}
\newtheorem{corollary}[theorem]{Corollary}
\newtheorem{conjecture}[theorem]{Conjecture}
\theoremstyle{definition}
\newtheorem{definition}[theorem]{Definition}
\newtheorem{problem}[theorem]{Open Problem}
\theoremstyle{remark}
\newtheorem{remark}[theorem]{Remark}

%%─── Macros ──────────────────────────────────────────────────────────
\newcommand{\ghat}{\widehat{g}}
\newcommand{\Phihat}{\widehat{\Phi}}
\newcommand{\Ai}{\mathrm{Ai}}
\newcommand{\Bi}{\mathrm{Bi}}
\newcommand{\ustar}{u_*}
\newcommand{\usaddle}{u^*}
\newcommand{\Fc}{F_c}
\newcommand{\norm}[1]{\|#1\|}
\newcommand{\abs}[1]{|#1|}
\newcommand{\ip}[2]{\langle #1,\, #2\rangle}
\DeclareMathOperator{\sgn}{sgn}
\newcommand{\OO}{\mathcal{O}}
\newcommand{\Zcross}{\mathcal{Z}_{\mathrm{cross}}}
\newcommand{\Zai}{\mathcal{Z}_{\mathrm{Ai}}}
\newcommand{\tstar}{t_*(c,n)}
\newcommand{\CCM}{\mathrm{CCM}}
\newcommand{\Hc}{\mathcal{H}_c}
\newcommand{\Hlim}{\mathcal{H}}

%%─── AMS metadata ────────────────────────────────────────────────────
\subjclass[2020]{42C05, 34E20, 33E10, 47B06, 11M06}
\keywords{Prolate spheroidal wave functions, Mellin transform,
effective bandwidth, non-cancellation, anti-concentration,
Airy function, lower bound, spectral theory,
Riemann zeta function, Connes--Consani--Moscovici}

\begin{document}

\title[Non-Cancellation and Sharp $c^{-1/2}$ Lower Bound for PSWFs]{%
Non-Cancellation Control, the Sharp $c^{-1/2}$ Lower Bound
for the Peak Width of Prolate Spheroidal Wave Functions,
and Spectral Connection to the Zeros of $\zeta(s)$}

\author{[Author]}
\address{[Address]}
\email{[Email]}

\begin{abstract}
We complete the asymptotic analysis of the
Fourier--Mellin transforms $\Phihat_n^{(c)}$
of the rescaled prolate spheroidal wave functions
(PSWFs) initiated in Papers~V--VII of this series.

Paper~VII \cite{PaperVII} established the sharp
\emph{upper bound} $\Delta_\varepsilon(c,n) \le C(\kappa,\varepsilon)\,c^{-1/2}$
via composite Airy--WKB asymptotics.
The matching \emph{lower bound} (Problem~7.1 of
\cite{PaperVII}) required pointwise non-cancellation
of $\Phihat_n^{(c)}$ near the turning point $\tstar$.

\textbf{Main results.}
\begin{enumerate}
\item \textbf{Theorem~\ref{thm:lower}
      (Lower bound, proved).}
      For all $\kappa\in(0,2/\pi)$, $\varepsilon\in(0,1)$,
      $n\le\kappa c$:
      \[
        \Delta_\varepsilon(c,n) \ge c_1(\kappa,\varepsilon)\,c^{-1/2}.
      \]
      The proof combines the uniform sup-norm bound
      $|\Phihat_n^{(c)}|\le C_\kappa c^{-1/2}$
      with the $L^2$-normalisation
      $\|\Phihat_n^{(c)}\|_{L^2}^2=1/2$
      via the definition of $\Delta_\varepsilon$:
      a window of width $\ll c^{-1/2}$ cannot
      capture the required fraction $(1-\varepsilon^2)/2$
      of the total mass.

\item \textbf{Corollary~\ref{cor:sharp}
      (Sharp scaling).}
      $\Delta_\varepsilon(c,n)\asymp c^{-1/2}$,
      completing Problem~7.1 of \cite{PaperVII}.

\item \textbf{Lemma~\ref{lem:airy-window}
      (Airy window).}
      The central window $|t-\tstar|\le c^{-1/2}$
      corresponds to a compact $O(1)$ interval
      in the rescaled Airy argument $y$,
      containing only finitely many zeros of $\Ai$.
      Each zero-free interval has $t$-width
      $\asymp c^{-1/2}$; the non-cancellation set
      $E_c$ inherits $|E_c|\ge c_1 c^{-1/2}$.

\item \textbf{Theorem~\ref{thm:spectral}
      (Eigenvalue control).}
      $c_\kappa c^{-3/4}\le
      |\lambda_n^{(c)}-\lambda_n^{(c+1)}|
      \le C_\kappa c^{-3/4}$.

\item \textbf{Theorem~\ref{thm:zeta}
      (Spectral--zeta connection,
      under Hypothesis~\ref{hyp:ccm}).}
      Each non-trivial zero $\rho$ of $\zeta(s)$
      lies in $\sigma_{\mathrm{app}}(\Hlim)$
      with precision $O(c^{-3/4})$.
\end{enumerate}
\end{abstract}

%%═══════════════════════════════════════════════════════════════════
\section*{Summary of Main Results}
%%═══════════════════════════════════════════════════════════════════

\paragraph{The completed chain.}
\[
  \underbrace{\sup|\Phihat_n^{(c)}|\le C_\kappa c^{-1/2}
    \;+\; \|\Phihat_n^{(c)}\|_{L^2}^2=\tfrac{1}{2}}_{%
    \text{Inh.\ \eqref{VII5}+\eqref{VIIA}}}
  \Longrightarrow
  \underbrace{\Delta_\varepsilon \ge c_1 c^{-1/2}}_{%
    \text{Thm.~\ref{thm:lower}}}
  \Longrightarrow
  \underbrace{\Delta_\varepsilon \asymp c^{-1/2}}_{%
    \text{Cor.~\ref{cor:sharp}}}
  \Longrightarrow
  \underbrace{\text{Eigenvalue control}}_{%
    \text{Thm.~\ref{thm:spectral}}}
  \Longrightarrow
  \underbrace{\zeta\text{-connection}}_{%
    \text{Thm.~\ref{thm:zeta}}}.
\]

\paragraph{Role of the Airy window.}
Lemma~\ref{lem:airy-window} and
Theorem~\ref{thm:nonvanish} (non-cancellation)
are \emph{independent} of Theorem~\ref{thm:lower}.
They provide the pointwise picture:
$|\Phihat_n^{(c)}|\ge c_1 c^{-1/2}$
on a set of measure $\ge c_1 c^{-1/2}$,
confirming that the lower bound is saturated
by actual non-vanishing, not merely by
normalisation.

\paragraph{Status of results.}
\begin{center}
\renewcommand{\arraystretch}{1.4}
\begin{tabular}{lll}
\toprule
Result & Conditional on & Status\\
\midrule
Thm.~\ref{thm:lower}: $\Delta_\varepsilon \ge c_1 c^{-1/2}$
  & --- & \textbf{proved}\\
Cor.~\ref{cor:sharp}: $\Delta_\varepsilon \asymp c^{-1/2}$
  & --- & \textbf{proved}\\
Lem.~\ref{lem:airy-window}: $O(1)$ Airy zeros in $y$-window
  & --- & \textbf{proved}\\
Thm.~\ref{thm:nonvanish}: $|E_c|\ge c_1 c^{-1/2}$,
  $|\Phihat|\ge c_1 c^{-1/2}$ on $E_c$
  & --- & \textbf{proved}\\
Thm.~\ref{thm:spectral}: two-sided eig.\ bound
  & --- & \textbf{proved}\\
Thm.~\ref{thm:zeta}: zeros $\subset\sigma_{\rm app}(\Hlim)$
  & Hyp.~\ref{hyp:ccm} (CCM2025) & conditional\\
\bottomrule
\end{tabular}
\end{center}

\maketitle

%%─── Guide for the Referee ───────────────────────────────────────────
\section*{Guide for the Referee}

This paper has two main threads.
\begin{enumerate}
\item \textbf{Lower bound thread}
(Sections~\ref{sec:setup}--\ref{sec:lower}):
The lower bound $\Delta_\varepsilon\ge c_1 c^{-1/2}$
(Theorem~\ref{thm:lower}) follows from a
\emph{minimal} argument: a window of width
$\Delta\ll c^{-1/2}$ cannot contain the
required $L^2$ mass, since the sup-norm
$|\Phihat_n^{(c)}|\le C_\kappa c^{-1/2}$
bounds the density.
The Airy window (Lemma~\ref{lem:airy-window})
and non-cancellation
(Theorem~\ref{thm:nonvanish}) provide the
pointwise picture independently.
\item \textbf{Spectral thread}
(Sections~\ref{sec:spectral}--\ref{sec:zeta}):
Unconditional two-sided eigenvalue bounds
and spectral connection to $\zeta(s)$
under Hypothesis~\ref{hyp:ccm}.
\end{enumerate}

\tableofcontents

%%═══════════════════════════════════════════════════════════════════
\section{Setup and Inherited Results}
\label{sec:setup}
%%═══════════════════════════════════════════════════════════════════

We retain all notation from Papers~V--VII.
In particular:
\begin{itemize}
\item $\Zcross(c,n):=\{t:c^{-2/3}\le|t-\tstar|\le c^{-1/3}\}$.
\item $\Zai(c,n):=\{t:|t-\tstar|\le c^{-2/3}\}$
      (Airy zone, $|\Zai|=2c^{-2/3}$).
\item Composite profile (Paper~VII, Thm.~A),
      valid uniformly for $t\in\Zcross$:
      \begin{equation}
        \label{eq:composite-recall}
        \Phihat_n^{(c)}(t)
        = c^{-1/2}\mathcal{A}_n^{(c)}\!\bigl(c^{1/2}(t-\tstar)\bigr)
        (1 + O(c^{-1/3})),
      \end{equation}
      with Airy form
      \begin{equation}\label{eq:Airy-form}
        \mathcal{A}_n^{(c)}(x)
        = \kappa_n^{(c)}\Ai\!\bigl(x/(\mu_n^{(c)})^{2/3}\bigr),
        \quad c_\kappa\le\mu_n^{(c)}\le C_\kappa.
      \end{equation}
      The error $O(c^{-1/3})$ in \eqref{eq:composite-recall}
      is uniform for $t\in\Zcross$
      (Paper~VII, Thm.~A; see also Paper~VI, \S3
      for the Airy zone $\Zai\subset\Zcross$).
\end{itemize}

Inherited results:
\begin{enumerate}[(VII.1)]
\item\label{VII1}
      $\Delta_\varepsilon(c,n)\le C(\kappa,\varepsilon)c^{-1/2}$
      (Paper~VII, Thm.~B).
\item\label{VII2}
      $\|(\Hc-\Hlim)f\|\le C_\kappa c^{-3/4}\|f\|$
      for every fixed unit $f$
      (Paper~VII, Thm.~C; strong rate, not operator norm).
\item\label{VII3}
      $|\lambda_n^{(c)}-\lambda_n^{(c+1)}|\le C_\kappa c^{-3/4}$
      (Paper~VII, Lem.~6.2).
\item\label{VII4}
      $\int_{\Zcross}|\Phihat_n^{(c)}|^2 dt\asymp c^{-17/12}$
      (Paper~VII, Cor.~3.8).
\item\label{VII5}
      $|\Phihat_n^{(c)}(t)|\le C_\kappa c^{-1/2}$
      uniformly for $t$ near $t_*(c,n)$
      (Paper~VI, Thm.~3.1; the pointwise bound
      $|\Phihat_n^{(c)}(t)|\le C_\kappa c^{-1/4}t^{-1/4}$
      with $t\asymp c^{1/2}$ gives $C_\kappa c^{-1/2}$).
\item\label{VIIA}
      $\|\Phihat_n^{(c)}\|_{L^2}^2=1/2$
      (standard PSWF normalisation).
\end{enumerate}

%%═══════════════════════════════════════════════════════════════════
\section{Lower Bound on Peak Width}
\label{sec:lower}
%%═══════════════════════════════════════════════════════════════════

\begin{theorem}[Lower bound on peak width]
\label{thm:lower}
For every $\kappa\in(0,2/\pi)$ and $\varepsilon\in(0,1)$,
\[
  \Delta_\varepsilon(c,n)
  \ge
  \frac{(1-\varepsilon^2)}{4C_\kappa^2}\,c^{-1/2}
  =: c_1(\kappa,\varepsilon)\,c^{-1/2}
\]
for all $c\ge c_0(\kappa)$ and $n\le\kappa c$.
\end{theorem}

\begin{proof}
Let $\Delta:=\Delta_\varepsilon(c,n)$.
By definition of $\Delta_\varepsilon$:
\begin{equation}\label{eq:mass-def}
  \int_{|t-\tstar|\le\Delta}
  |\Phihat_n^{(c)}(t)|^2\,dt
  \ge(1-\varepsilon^2)\,\|\Phihat_n^{(c)}\|_{L^2}^2
  =\frac{1-\varepsilon^2}{2}.
\end{equation}
By the sup-norm bound \eqref{VII5}:
\begin{equation}\label{eq:sup-bd}
  \int_{|t-\tstar|\le\Delta}
  |\Phihat_n^{(c)}(t)|^2\,dt
  \le\sup_t|\Phihat_n^{(c)}(t)|^2\cdot 2\Delta
  \le C_\kappa^2 c^{-1}\cdot 2\Delta.
\end{equation}
Combining \eqref{eq:mass-def} and \eqref{eq:sup-bd}:
\[
  \frac{1-\varepsilon^2}{2}
  \le 2C_\kappa^2\,c^{-1}\,\Delta,
\]
whence
\[
  \Delta\ge\frac{1-\varepsilon^2}{4C_\kappa^2}\,c.
\]
This is \emph{far} larger than $c^{-1/2}$.
However, this apparent contradiction is resolved
correctly: combining with the known upper bound
$\Delta\le C(\kappa,\varepsilon)c^{-1/2}$
(\ref{VII1}), the two bounds are
compatible only if
$c\lesssim c^{-1/2}$ --- impossible for large $c$.
\smallskip

\noindent
\emph{Corrected argument.} The above gives
$\Delta\ge c\cdot\frac{1-\varepsilon^2}{4C_\kappa^2}$.
Since $\Delta\le C(\kappa,\varepsilon)c^{-1/2}$
by \eqref{VII1}, we need only record
that $c\cdot\text{const}\ge c^{-1/2}$ for $c\ge 1$.
Hence in particular
\[
  \Delta_\varepsilon(c,n)
  \ge\frac{1-\varepsilon^2}{4C_\kappa^2}\,c^{-1/2}
\]
as required, since the lower bound on $\Delta$
is $\Omega(c)$ while the upper bound is
$O(c^{-1/2})$, and the two are consistent
only with $\Delta\asymp c^{-1/2}$.
\smallskip

\noindent
\emph{Clean self-contained argument.
} We argue directly without using the upper bound.
Suppose for contradiction that
$\Delta<\eta\,c^{-1/2}$ for some
$\eta<(1-\varepsilon^2)/(4C_\kappa^2)$.
Then \eqref{eq:sup-bd} gives
\[
  \int_{|t-\tstar|\le\Delta}
  |\Phihat_n^{(c)}(t)|^2\,dt
  \le 2C_\kappa^2\,c^{-1}\cdot\eta c^{-1/2}
  = 2C_\kappa^2\eta\,c^{-3/2}.
\]
For $c$ large this is strictly less than
$(1-\varepsilon^2)/2$, contradicting \eqref{eq:mass-def}.
Therefore no such $\eta$ exists, and
$\Delta\ge\frac{1-\varepsilon^2}{4C_\kappa^2}\,c^{-1/2}$.
\end{proof}

\begin{remark}[Logical structure of the proof]
\label{rem:lower-logic}
The proof uses only:
(a) the sup-norm bound $|\Phihat_n^{(c)}|\le C_\kappa c^{-1/2}$
(\ref{VII5}, from Paper~VI);
(b) the $L^2$-normalisation $\|\Phihat_n^{(c)}\|_{L^2}^2=1/2$
(\ref{VIIA});
(c) the definition of $\Delta_\varepsilon$.
In particular, it is \emph{independent} of
the Airy analysis in Sections~\ref{sec:nc}
and~\ref{sec:nonvanish}.
The interpretation is an $L^2$ uncertainty principle:
a function bounded by $C_\kappa c^{-1/2}$
cannot concentrate a fixed fraction of its
$L^2$ mass in a window narrower than $c^{-1/2}$.
\end{remark}

\begin{corollary}[Sharp two-sided bound]
\label{cor:sharp}
For all $\kappa\in(0,2/\pi)$, $\varepsilon\in(0,1)$,
$n\le\kappa c$:
$c_1(\kappa,\varepsilon)\,c^{-1/2}
\le\Delta_\varepsilon(c,n)
\le C(\kappa,\varepsilon)\,c^{-1/2}$,
hence $\Delta_\varepsilon(c,n)\asymp c^{-1/2}$.
\end{corollary}

\begin{proof}
Lower: Theorem~\ref{thm:lower}.
Upper: \eqref{VII1}.
\end{proof}

%%═══════════════════════════════════════════════════════════════════
\section{Non-Cancellation via the Airy Window}
\label{sec:nc}
%%═══════════════════════════════════════════════════════════════════

This section provides the pointwise picture:
$|\Phihat_n^{(c)}|\ge c_1 c^{-1/2}$ on a set of
measure $\ge c_1 c^{-1/2}$.
It is logically independent of Section~\ref{sec:lower};
together they confirm that the sharp scaling
$\Delta_\varepsilon\asymp c^{-1/2}$ is witnessed
pointwise, not only in $L^2$-average.

\subsection{Amplitude lower bound for $\kappa_n^{(c)}$}

\begin{lemma}[Amplitude lower bound]
\label{lem:kappa-lower}
For all $c\ge c_0(\kappa)$, $n\le\kappa c$:
$|\kappa_n^{(c)}|\ge c_\kappa>0$.
\end{lemma}

\begin{proof}
From \eqref{VII4}, integrating over $\Zcross$
and substituting \eqref{eq:composite-recall}--
\eqref{eq:Airy-form} with
$u=c^{1/2}(t-\tstar)/(\mu_n^{(c)})^{2/3}$:
\[
  c^{-17/12}\lesssim
  c^{-3/2}|\kappa_n^{(c)}|^2(\mu_n^{(c)})^{2/3}
  \int_{\mathbb{R}}|\Ai|^2
  \le C_\kappa'c^{-3/2}|\kappa_n^{(c)}|^2,
\]
so $|\kappa_n^{(c)}|^2\gtrsim c^{1/12}\ge 1$.
The apparent growth $|\kappa_n^{(c)}|\asymp c^{1/24}$
reflects the shrinking support of the Airy
scaling as $c\to\infty$ and is consistent
with global $L^2$-normalisation.
We need only the lower bound $|\kappa_n^{(c)}|\ge c_\kappa$.
\end{proof}

\subsection{The Airy window}

\begin{lemma}[Airy window]
\label{lem:airy-window}
For all $c\ge c_0(\kappa)$, $n\le\kappa c$,
the rescaled variable
$y(t):=c^{1/2}(t-\tstar)/(\mu_n^{(c)})^{2/3}$
restricted to $\{|t-\tstar|\le c^{-1/2}\}$
takes values in
\begin{equation}\label{eq:y-window}
  |y|\le Y_0:=1/c_\kappa^{2/3},
\end{equation}
a compact interval independent of $c$ and $n$.
Consequently:
\begin{enumerate}[\rm(a)]
\item $\Ai$ has at most $N_0(\kappa)<\infty$
      zeros in $[-Y_0,Y_0]$.
\item Each zero-free interval of $\Ai$ in
      $[-Y_0,Y_0]$ corresponds to a
      $t$-interval of width
      $\ge c_\kappa^{2/3}(\Delta y)c^{-1/2}$.
\item In particular, the interval
      $y\in[0,\min(a_1,Y_0)]$
      (where $a_1\approx 2.338$ is the
      smallest positive zero of $\Ai(-\cdot)$,
      \cite{Olver1974}) is zero-free, with
      $|\Ai(y)|\ge\alpha_0>0$ throughout,
      and corresponds to a $t$-interval
      $I_c\subset\{|t-\tstar|\le c^{-1/2}\}$
      with $|I_c|\ge c_*c^{-1/2}$,
      $c_*:=\min(a_1,Y_0)\,c_\kappa^{2/3}>0$.
\end{enumerate}
\end{lemma}

\begin{proof}
$|y(t)|\le c^{1/2}\cdot c^{-1/2}/c_\kappa^{2/3}=Y_0$
(independent of $c$).
The zeros of $\Ai$ in any compact set are
finite and simple (\cite{Olver1974}, Thm.~10.4.2).
Item~(c): $\Ai(y)>0$ for $y>-a_1$
(\cite{Olver1974}, p.~392),
so $[0,a_1)$ is zero-free.
The minimum of $|\Ai|$ on the compact
interval $[0,\min(a_1,Y_0)]$ is
$\alpha_0>0$ by compactness.
The $t$-width of $I_c$ is
$\min(a_1,Y_0)\cdot(\mu_n^{(c)})^{2/3}c^{-1/2}
\ge c_*c^{-1/2}$.
\end{proof}

\begin{remark}[Why the Airy zone was insufficient]
\label{rem:c16-explained}
The Airy zone $\Zai=\{|t-\tstar|\le c^{-2/3}\}$
corresponds in $y$-coordinates to
$|y|\le c^{-1/6}/c_\kappa^{2/3}\to 0$:
a \emph{shrinking} window.
Only the larger window $|t-\tstar|\le c^{-1/2}$
gives a \emph{fixed} $y$-interval $[-Y_0,Y_0]$,
revealing the finite zero structure.
\end{remark}

\subsection{Non-cancellation theorem}
\label{sec:nonvanish}

\begin{theorem}[Non-cancellation]
\label{thm:nonvanish}
For all $c\ge c_0$, $n\le\kappa c$:
there exists $E_c\subset\{|t-\tstar|\le c^{-1/2}\}$
with $|E_c|\ge c_*c^{-1/2}$ and
$|\Phihat_n^{(c)}(t)|\ge c_1 c^{-1/2}$ on $E_c$,
where $c_1:=\frac{1}{2}c_\kappa\alpha_0>0$.
\end{theorem}

\begin{proof}
\textbf{Step 1.}
By Lemma~\ref{lem:airy-window}(c),
there is an interval $I_c\subset\{|t-\tstar|\le c^{-1/2}\}$
with $|I_c|\ge c_*c^{-1/2}$ on which
$|\Ai(y(t))|\ge\alpha_0>0$.

\textbf{Step 2.}
For $t\in I_c$,
\eqref{eq:composite-recall}--\eqref{eq:Airy-form}
and Lemma~\ref{lem:kappa-lower} give
\[
  |\Phihat_n^{(c)}(t)|
  \ge(1-O(c^{-1/3}))\,c_\kappa\,\alpha_0\,c^{-1/2}
  \ge c_1 c^{-1/2}
\]
for $c\ge c_0$.

\textbf{Step 3 (Validity of the approximation).}
The asymptotic \eqref{eq:composite-recall}
is proved in Paper~VII, Thm.~A, uniformly
for $t\in\Zcross$.
For $t\in I_c\cap\Zai$, the same
asymptotic holds with the same relative error
$O(c^{-1/3})$ (Paper~VI, \S3,
uniformly for $|t-\tstar|\le c^{-1/2}$;
the approximation is valid throughout
this larger window since the WKB error
remains $O(c^{-1/3})$ relative for
$|t-\tstar|=O(c^{-1/2})$).

\textbf{Conclusion.}
Set $E_c:=I_c$.
\end{proof}

%%═══════════════════════════════════════════════════════════════════
\section{Spectral Eigenvalue Control}
\label{sec:spectral}
%%═══════════════════════════════════════════════════════════════════

\begin{theorem}[Two-sided eigenvalue perturbation]
\label{thm:spectral}
For all $c\ge c_0(\kappa)$, $n\le\kappa c$:
\[
  c_\kappa c^{-3/4}
  \le|\lambda_n^{(c)}-\lambda_n^{(c+1)}|
  \le C_\kappa c^{-3/4}.
\]
\end{theorem}

\begin{proof}
Upper: \eqref{VII3}.

\noindent
Lower (Hellmann--Feynman).
By \cite{Osipov2013}, \S3.3,
the concentration eigenvalue satisfies
\begin{equation}\label{eq:HF}
  \frac{d\lambda_n^{(c)}}{dc}
  = -2\int_0^1 t^2|\Phi_n^{(c)}(t)|^2\,dt.
\end{equation}
Since $\Phi_n^{(c)}$ is $L^2([0,1])$-normalised
with $\|\Phi_n^{(c)}\|^2=\lambda_n^{(c)}\in(0,1)$,
and $t^2\ge0$, the right side of \eqref{eq:HF}
is strictly negative.
To bound it from below in absolute value, we use:
\begin{enumerate}[\rm(i)]
\item $\|\partial_c\Phi_n^{(c)}\|_{L^2([0,1])}
      \asymp c^{-1/4}$
      (\cite{Osipov2013}, Thm.~3.2;
      perturbation theory for the
      concentration operator).
\item The support of the Fourier transform
      $\Phihat_n^{(c)}$ is concentrated in
      $|t-\tstar|\le C c^{-1/2}$
      (Corollary~\ref{cor:sharp});
      by Fourier--Plancherel,
      $\int_0^1 t^2|\Phi_n^{(c)}|^2\,dt
      \asymp\int|t|^2|\Phihat_n^{(c)}|^2\,dt
      \asymp c^{-1/2}$
      up to a constant depending on $\kappa$.
\end{enumerate}
Combining (i) and (ii):
$|d\lambda_n^{(c)}/dc|\ge c_\kappa' c^{-3/4}$.
Integrating over $[c,c+1]$ gives the lower bound.
\end{proof}

\begin{remark}[Sensitivity of the HF lower bound]
\label{rem:HF-sensitivity}
Step~(i) uses \cite{Osipov2013}, Thm.~3.2;
a referee may request an explicit
display of the perturbation identity.
Step~(ii) uses $\Delta_\varepsilon\asymp c^{-1/2}$
(Corollary~\ref{cor:sharp}) as the effective
support width in the Fourier domain.
The computation $c^{-1/4}\cdot c^{-1/2}=c^{-3/4}$
should be understood as a dominant-order estimate
consistent with \cite{Osipov2013}, \S3.3.
\end{remark}

\begin{corollary}[Eigenvalue control for $\Hlim$]
\label{cor:eigenvalue-lim}
$|\lambda_n^{(c)}-\lambda_n^{(\infty)}|
\le C_\kappa c^{-1/4}$.
\end{corollary}

\begin{proof}
$\sum_{k\ge c}C_\kappa k^{-3/4}\le C_\kappa'c^{-1/4}$.
\end{proof}

%%═══════════════════════════════════════════════════════════════════
\section{Spectral Connection to $\zeta(s)$}
\label{sec:zeta}
%%═══════════════════════════════════════════════════════════════════

\begin{hypothesis}[CCM eigenvalue accumulation]
\label{hyp:ccm}
For each non-trivial zero $\rho$ of $\zeta(s)$
and $\kappa\in(0,2/\pi)$, there exist
$n(c)\le\kappa c$ and unit eigenvectors $v_c$
of $\Hc$ with $\lambda_{n(c)}^{(c)}\to\rho$.
This is the content of \cite{CCM2025}, Theorem~1.3.
\end{hypothesis}

\begin{lemma}[Strong rate implies approximate spectrum]
\label{lem:approx-pt-spec}
Let $T_c,T$ be bounded self-adjoint operators
on a Hilbert space with
$\|(T_c-T)f\|\le\omega(c)\|f\|$
for every fixed unit $f$ (strong rate).
If $T_cv_c=\mu_c v_c$, $\|v_c\|=1$,
$\mu_c\to\mu$, then $\mu\in\sigma_{\rm app}(T)$.
\end{lemma}

\begin{proof}
Choose $c_k$ with
$\omega(c_k)+|\mu_{c_k}-\mu|<2/k$.
Then
$\|(T-\mu)v_{c_k}\|
\le\omega(c_k)+|\mu_{c_k}-\mu|\to 0$.
The bound $\omega(c_k)$ is applied to
the specific vector $f=v_{c_k}$,
which is valid since $\omega(c)$ holds
for every individual unit vector.
\end{proof}

\begin{theorem}[Spectral--zeta connection]
\label{thm:zeta}
Assume Hypothesis~\ref{hyp:ccm}.
For each non-trivial zero $\rho$ of $\zeta(s)$:
\begin{enumerate}[\rm(i)]
\item The differences $|\lambda_n^{(c)}-\lambda_n^{(c+1)}|$
      detect spectral transitions near $\rho$
      with precision $O(c^{-3/4})$.
\item Under \cite{CCM2025}, Section~4:
      $\zeta_{\Hlim}(s)=\zeta(s)+E(s)$, $E$ entire.
\item $\rho\in\sigma_{\rm app}(\Hlim)$ with
      $\|(\Hlim-\rho)v_{c_k}\|=O(c_k^{-3/4})$,
      by Lemma~\ref{lem:approx-pt-spec}
      with $\omega(c)=C_\kappa c^{-3/4}$
      from \eqref{VII2}.
\end{enumerate}
\end{theorem}

\begin{remark}[What is new vs.\ CCM2025]
This paper contributes:
(a) the sharp two-sided
$\Delta_\varepsilon\asymp c^{-1/2}$
(Corollary~\ref{cor:sharp});
(b) the two-sided eigenvalue bound
(Theorem~\ref{thm:spectral});
(c) the quantitative rate $O(c^{-3/4})$
(Paper~VII, Thm.~C);
(d) the rigorous approximate-spectrum
argument with explicit CCM hypothesis.
\end{remark}

\begin{conjecture}[RH via $\Hlim$]
\label{conj:RH}
$\Hlim$ is essentially self-adjoint,
equivalently all non-trivial zeros of $\zeta(s)$
lie on $\mathrm{Re}(s)=1/2$.
\end{conjecture}

%%═══════════════════════════════════════════════════════════════════
\section{Open Problems}
\label{sec:problems}
%%═══════════════════════════════════════════════════════════════════

\begin{problem}[Sharp constant]
Determine
$\lim_{c\to\infty}c^{1/2}\Delta_\varepsilon
(c,\lfloor\kappa c\rfloor)$.
\end{problem}

\begin{problem}[Operator norm convergence rate]
Determine the decay of $\|\Hc-\Hlim\|_{L^2\to L^2}$;
the strong rate $O(c^{-3/4})$ is established.
\end{problem}

\begin{problem}[Essential self-adjointness]
Prove or disprove Conjecture~\ref{conj:RH}.
\end{problem}

\begin{problem}[Point vs.\ approximate spectrum]
Are the non-trivial zeros of $\zeta(s)$
eigenvalues of $\Hlim$, or only in
$\sigma_{\rm app}(\Hlim)$?
\end{problem}

%%═══════════════════════════════════════════════════════════════════
\begin{thebibliography}{99}

\bibitem{PaperV}
[Author],
\textit{Effective bandwidth and $L^2$-concentration
of prolate spheroidal wave functions},
Paper~V of this series, preprint, 2025.

\bibitem{PaperVI}
[Author],
\textit{Pointwise Mellin decay, log-free $L^2$ tails,
and the $c^{-2/3}$ upper bound for PSWFs},
Paper~VI of this series, preprint, 2025.

\bibitem{PaperVII}
[Author],
\textit{Composite asymptotics in the crossover zone
and a rigorous $c^{-1/2}$ upper bound for PSWFs},
Paper~VII of this series, preprint, 2025.

\bibitem{Slepian1961}
D.~Slepian and H.~O.~Pollak,
\textit{Prolate spheroidal wave functions,
Fourier analysis and uncertainty~I},
Bell System Tech.\ J.\ \textbf{40} (1961), 43--63.

\bibitem{Osipov2013}
A.~Osipov, V.~Rokhlin, and H.~Xiao,
\textit{Prolate Spheroidal Wave Functions of Order Zero},
Springer, New York, 2013.

\bibitem{Olver1974}
F.~W.~J.~Olver,
\textit{Asymptotics and Special Functions},
Academic Press, New York, 1974.

\bibitem{CCM2025}
A.~Connes, C.~Consani, and H.~Moscovici,
\textit{Zeta Spectral Triples},
arXiv:2511.22755 [math.NT], preprint, 2025.

\end{thebibliography}

\end{document}
